\chapter{Impostazione sperimentale}
\label{sec:setup}

% Source notes: Project/docs/thesis_notes.md; checkpoint_2026-10-10.md; configs Project/GroundOne/configs/train_0.5b.yaml,
% train_3b.yaml; queues Project/jobs/queue_night3.txt, queue_day4.txt, queue_night5.txt, queue_day5.txt, queue_night6.txt.

Questo capitolo presenta i modelli addestrati, i sistemi con cui sono messi a confronto, le metriche, il protocollo
statistico e le modalità di misura di tempo ed energia. Tutti gli esperimenti si svolgono su un'unica macchina, un
MacBook Pro con chip Apple M4 (CPU a 10 core, GPU a 10 core, banda di memoria di 120\,GB/s) e 16\,GB di memoria
unificata~\cite{apple2024m4specs}: ogni risultato riflette così ciò che un verificatore può fare sullo stesso dispositivo
su cui gira un assistente locale. Anche tutte le baseline, compresi gli encoder e i due System One generalisti, sono
eseguite su questo portatile, e nessun risultato dipende da un servizio remoto.

\section{Modelli}
\label{sec:setup-models}

Entrambi i modelli GroundOne partono da Qwen2.5~\cite{qwen2024qwen25} nella versione instruction-tuned, convertita in
MLX con pesi a 8 bit, e sono addestrati con la procedura descritta nella Sezione~\ref{sec:method-training}. Si
distinguono soltanto per la dimensione, per il numero di risposte di addestramento e per la composizione dei
micro-batch, che dipende dalla memoria disponibile.

Il modello 0.5B è addestrato sull'intero training set di 13{,}203 risposte. Un'epoca di 802 passi di ottimizzazione ha
richiesto 5.5 ore di tempo wall-clock, con una valutazione sull'intero dev set ogni 150 passi, e un'occupazione massima
di 2.6\,GB di memoria. Il modello 3B è circa sei volte più lento per token; per questo viene addestrato per una notte su
un sottoinsieme casuale di 6{,}500 risposte (circa 45{,}000 decisioni), con 403 passi in 12.25 ore e un picco di
5.9\,GB. La sua valutazione sul dev set, ogni 80 passi, usa un sottoinsieme fisso di 200 risposte, dato che una
valutazione completa richiederebbe circa un'ora. In entrambi i casi si conserva il checkpoint con il miglior punteggio
sul dev set, che viene poi calibrato come descritto nella Sezione~\ref{sec:method-calibration}; per il 3B si tratta del
passo 320.

Per stimare quanto i risultati dipendano dal seed, che determina l'inizializzazione degli adapter e l'ordine dei batch,
la procedura del 0.5B viene ripetuta con un secondo seed: è noto che l'esito del fine-tuning varia con entrambi questi
fattori~\cite{dodge2020finetuning}. Un'ulteriore esecuzione con un terzo seed è collassata presto su una predizione
costante ed è stata interrotta a circa il 60\,\% dell'epoca; poiché incompleta, è esclusa dalle tabelle e viene discussa
nella Sezione~\ref{sec:results-seeds}. La seconda esecuzione, portata a termine, è valutata su tutti e tre i test set
esattamente come la prima.

\section{Baseline}
\label{sec:setup-baselines}

Le baseline servono a isolare l'origine di un eventuale vantaggio, e ciascuna risponde a una domanda precisa.

\paragraph{Segnali che non leggono mai il documento nella sua interezza.} Due predittori elementari fissano la soglia
minima che qualsiasi verificatore deve superare. Il primo assegna a una frase un punteggio pari alla frazione delle sue
parole di contenuto che compaiono anche nel documento. Il secondo è un classificatore bag-of-words addestrato sulle sole
frasi del training set, senza il documento. Entrambi sono descritti nella Sezione~\ref{sec:data-eda}.

\paragraph{Readout senza addestramento (B1).} Agli stessi modelli Qwen2.5, senza fine-tuning, si pone per ogni frase la
domanda a tre classi e se ne leggono i logit delle lettere di risposta. Un modello non addestrato non ha alcun motivo di
interpretare il marker \texttt{Verdict:} come una domanda, perciò questa baseline adotta il formato con una domanda
esplicita per frase (Sezione~\ref{sec:method-readout}). Misura quindi ciò che il meccanismo System One offre prima di
qualunque specializzazione.

\paragraph{Giudici generativi (B0a, B0b).} I modelli non addestrati sono impiegati anche nel modo più diffuso nella
pratica: ricevono il documento, la risposta numerata e l'elenco delle frasi da verificare, e scrivono un verdetto per
ogni frase con greedy decoding. Il giudice \emph{sintetico} produce una riga ``\texttt{<number>: <letter>}'' per frase;
il giudice \emph{verboso} scrive una breve frase di ragionamento seguita dal verdetto. Un verdetto che non è possibile
interpretare costituisce un errore di formato: viene contato come errore nel calcolo dell'accuratezza e riportato a
parte.

\paragraph{Giudici a parità di output (B0d).} I giudici sintetico e verboso restituiscono meno di GroundOne, cioè una
lettera per frase e nessuna probabilità, e non consentono quindi né calibrazione, né una soglia regolabile, né la scelta
di una politica di decisione. Per questo a due ulteriori giudici si chiede esattamente ciò che produce GroundOne:
un'etichetta e tre probabilità per frase. Il primo le scrive in un oggetto JSON, con le probabilità espresse a parole,
che è quanto uno sviluppatore otterrebbe con il solo prompt. Il secondo è il giudice sintetico, con le tre probabilità
lette dai logit in corrispondenza della lettera che genera; ottiene lo stesso contenuto al costo del giudice sintetico ed
è l'avversario generativo più forte. Entrambi sono calibrati con la stessa temperatura e lo stesso bias di GroundOne,
stimati su 200 risposte del dev set, così che le loro probabilità ricevano lo stesso trattamento. Il prompt JSON è stato
scelto su dodici risposte del dev set badando unicamente alla validità del formato: i modelli copiavano alla lettera un
esempio numerico presente nel prompt, che è stato perciò sostituito da segnaposto. Il parsing è permissivo, a vantaggio
dei giudici: le voci complete di un oggetto troncato vengono conservate e un'etichetta mancante si ricava dalle
probabilità. Il decoding non è vincolato da una grammatica~\cite{willard2023outlines}, in modo che il protocollo resti
quello degli altri giudici; questa scelta pesa soprattutto sugli errori di formato del modello più piccolo.

\paragraph{Readout e generazione a confronto (B2).} I modelli GroundOne addestrati sono eseguiti con lo stesso prompt e
lo stesso singolo forward pass, ma il verdetto a ogni marker è il token con la probabilità più alta sull'intero
vocabolario, cioè quello che produrrebbe la generazione greedy di un solo token. Le tre lettere, con o senza spazio
iniziale, corrispondono alle tre classi, e qualunque altro token è un errore di formato. Questo controllo lascia
invariati pesi, prompt e costo, e cambia soltanto se la decisione sia ristretta o meno alle opzioni ammissibili.

\paragraph{Un System One generalista: Rizzo-Flow (B3).} Il progetto open Rizzo-Flow~\cite{rizzoflow} realizza l'idea di System One
con un modello da 4B sottoposto a fine-tuning su decisioni tipizzate generiche. Gira in un server dedicato, basato su
llama.cpp~\cite{llamacpp}, e viene interrogato via HTTP passando il documento e la risposta numerata come stato
condiviso e una domanda a scelta fra tre opzioni per ogni frase. Si usa il modello quantizzato a 4 bit perché quello a
8 bit non entra in memoria insieme al sistema operativo.

\paragraph{Un System One generalista della stessa dimensione: d1-3B (B3b).} d1-3B di Liquid AI~\cite{liquid2026d1}
è un decision model open da 3.1B parametri, post-addestrato per decisioni tipizzate su un'ampia varietà di task, che legge
lettere di risposta ristrette con lo stato condiviso tra le domande: appartiene alla stessa famiglia di readout di
GroundOne e ha la dimensione del suo modello più grande, ma è generalista. Consente così di isolare il valore della
specializzazione rispetto a un System One già addestrato a decidere, anziché rispetto a un modello non addestrato.
Viene eseguito zero-shot con il proprio codice di inferenza, in PyTorch sulla GPU del portatile e senza accesso alla
rete, con lo stesso stato e le stesse domande di Rizzo-Flow e al più sei domande per chiamata sullo stesso stato; su
LLM-AggreFact i documenti sono suddivisi esattamente come per GroundOne e si conserva il massimo sui frammenti. Ai fini
della riproducibilità conta un dettaglio pratico: sulla GPU Apple PyTorch mantiene i kernel compilati per ogni nuova
lunghezza di input senza mai rilasciarne la memoria, perciò la valutazione si svolge in processi worker che vengono
riavviati quando la memoria occupata dal driver della GPU supera 11\,GB, senza alcun effetto sugli output.

\paragraph{Encoder specializzati (B4).} MiniCheck-RoBERTa~\cite{tang2024minicheck} e HHEM-2.1-open~\cite{vectara2024hhem}
sono i verificatori di piccole dimensioni più forti oggi disponibili. Valutano ogni coppia (documento, frase) in modo
indipendente e restituiscono un'unica probabilità di supporto; i documenti lunghi vengono suddivisi in frammenti di circa
400 parole e si conserva il massimo sui frammenti, come in MiniCheck. Poiché non hanno un output a tre classi, sono
confrontati soltanto sulla decisione binaria. Girano in PyTorch sulla GPU del portatile.

\paragraph{Risultati pubblicati.} Per LLM-AggreFact, i punteggi di sistemi più grandi presenti nella leaderboard
pubblica, come Qwen2.5-7B usato zero-shot, sono riportati così come pubblicati, senza rieseguirli.

I sistemi più lenti sono valutati su sottocampioni stratificati dei test set: Rizzo-Flow su 300 risposte di T1, 600
coppie di T3 e 1{,}100 elementi di T2, di cui ne ha valutati 1{,}093 (sette eccedevano il suo contesto); d1-3B su un
sottocampione casuale di 4{,}000 elementi di T2; il 3B non addestrato su 600 risposte di T1; i giudici generativi sulle
200 risposte del protocollo di latenza descritto più avanti. GroundOne, il 0.5B non addestrato, gli encoder e, su T1 e
T3, d1-3B sono valutati sui test set completi, e ogni confronto con un sistema sottocampionato è condotto sugli elementi
valutati da entrambi i sistemi.

\section{Metriche}
\label{sec:setup-metrics}

Poiché la quota di frasi allucinate non è la stessa nel dev set e in T1 (Sezione~\ref{sec:data-eda}), le metriche
primarie sono scelte in modo da non dipendere dal tasso di base. Per la decisione a tre classi la metrica principale è
la balanced accuracy, cioè la recall media sulle tre classi, affiancata dalla recall di ciascuna classe e dalla
frazione di predizioni per classe, che rivela un modello collassato su un'unica risposta. Per la decisione binaria si
riportano la balanced accuracy con soglia 0.5, come nella leaderboard di LLM-AggreFact, e l'AUROC di
$P(\text{supported})$, che non dipende da alcuna soglia. La calibrazione è misurata con l'ECE definito nella
Sezione~\ref{sec:method-calibration}, con il Brier score e con la negative log-likelihood; la selective prediction con
le curve risk--coverage e l'area sottesa (AURC)~\cite{geifman2017selective}. Tutte le metriche su T1 sono riportate
anche per task e per fascia di $K$.

\section{Protocollo statistico}
\label{sec:setup-statistics}

Ogni differenza tra due sistemi è accompagnata da un intervallo di confidenza al 95\,\% ottenuto con un paired bootstrap
con 2{,}000 resampling~\cite{koehn2004significance,bergkirkpatrick2012significance,dror2018hitchhiker}. L'unità di
resampling non è la frase, ma il gruppo impiegato nel controllo di leakage: il documento sorgente in RAGTruth, il
documento in LLM-AggreFact e la pagina di Wikipedia in VitaminC. Le frasi di una stessa risposta, e le risposte relative
a uno stesso documento, sono fortemente correlate, e ricampionarle in modo indipendente porterebbe a intervalli troppo
stretti. Su LLM-AggreFact i gruppi sono ricampionati all'interno di ciascuno degli undici dataset, perché la media macro
assegna a ogni dataset lo stesso peso. I due sistemi di un confronto sono valutati sugli stessi resampling, e il
$p$-value riportato è la frazione bilaterale di resampling in cui la differenza cambia segno.

Due limiti del protocollo vanno dichiarati in anticipo. Innanzitutto, si eseguono circa cinquanta confronti senza alcuna
correzione per confronti multipli: le conclusioni si fondano perciò su intervalli ben lontani dallo zero, e nessuna
differenza con un $p$-value prossimo a 0.05 viene considerata acquisita. Inoltre, il resampling dei dati non coglie la
variabilità dovuta all'addestramento, che è invece documentata a parte dalle due esecuzioni del 0.5B
(Sezione~\ref{sec:results-seeds}).

\section{Latenza, energia e hardware}
\label{sec:setup-protocol}

La latenza è misurata su 200 risposte di T1, 50 per ciascuna fascia di $K$, in modo che le risposte lunghe siano
rappresentate quanto quelle brevi. GroundOne, il readout con una domanda per frase e tutti i giudici generativi girano
in-process su MLX, una risposta alla volta, come farebbe un verificatore integrato in un assistente locale. Le prime tre
risposte fungono da warm-up e sono escluse. Per le varianti System One ogni risposta è cronometrata tre volte e si
conserva la mediana; i giudici generativi sono cronometrati una sola volta, dato che il greedy decoding è deterministico
e ogni esecuzione dura molto di più. Il tempo riportato è il tempo wall-clock della chiamata al modello per una
risposta, da cui si ricava il tempo per decisione.

Sono stati registrati i tempi anche per gli encoder, per Rizzo-Flow e per d1-3B, ma in condizioni che non permettono un
confronto equo. Gli encoder giravano in PyTorch e valutavano ogni frase separatamente, rileggendo ogni volta il
documento, il che li penalizza; Rizzo-Flow girava in un altro motore di inferenza, con pesi a 4 bit e dietro
un'interfaccia HTTP; d1-3B girava in PyTorch con il proprio codice e con i riavvii dei worker descritti sopra. Nei loro
confronti, quindi, non si avanza alcuna conclusione sulla latenza.

L'energia è misurata con \texttt{powermetrics}, lo strumento di macOS che riporta ogni 200 millisecondi la potenza
assorbita da CPU, GPU e Neural Engine. Per ciascun sistema, 100 risposte del sottoinsieme di latenza (25 per fascia di
$K$) vengono verificate una sola volta, dopo 60 secondi di misura a riposo con il modello già caricato; l'energia di una
verifica è la potenza media durante l'esecuzione, al netto della potenza a riposo, moltiplicata per la durata
dell'esecuzione e divisa per il numero di risposte. La misura comprende le unità di calcolo del chip ma non la memoria,
lo schermo o l'archiviazione, ed è stata effettuata di notte, con la macchina altrimenti inattiva.

È in esecuzione un solo modello alla volta, e una guardia di memoria sorveglia ogni job lungo, registrando per tutta la
sua durata la memoria libera e lo swap. I risultati vengono scritti una sola volta e mai sovrascritti, e ogni
esecuzione conserva la propria configurazione: ogni numero del Capitolo~\ref{sec:results} può così essere ricondotto
all'esecuzione che lo ha prodotto.
